\section{Energia cinética}
\label{sec_dinAcoplada_energia_cinetica}

A energia cinética total é escrita como a soma da contribuição associada à atitude da espaçonave (tratada como corpo rígido) e das contribuições de cada um dos $n$ elos do manipulador robótico, isto é:


\begin{equation}
T = T_{B} + \sum_{i=1}^{n} T_{i}.
\label{eq_T_total}
\end{equation}

O termo $n$ é a quantidade de elos do \gls{mr}, e $T_B$ é a energia cinética da espaçonave, em torno de seu \gls{cdm}, incluindo o \gls{mr} o modo de operação.
A Equação~\eqref{eq_T_total} poderia ser reescita com uma parcela linear, associada ao movimento do seu \gls{cdm}, e uma parcela rotacional, associada ao movimento de atitude em torno desse mesmo ponto.
Essa energia cinética de atitude é escrita como:

\begin{equation}
T_{B} =
\frac{1}{2}\, {}^{B}\!\vec{\omega}_{B}^{\,T}\,\mathbf{I}_B(\vec{\xi})\,{}^{B}\!\vec{\omega}_{B}.
\label{eq_TB}
\end{equation}

Em que $\mathbf{I}_B(\vec{\xi})$ é a matriz de inércia do perseguidor em torno do seu \gls{cdm} $G$, expressa no referencial $\{B\}$ fixo ao perseguidor.
Em particular, $\mathbf{I}_B(\vec{\xi})$ depende da configuração articular $\vec{\xi}$, uma vez que a redistribuição de massa provocada pelo movimento do \gls{mr} altera a posição do \gls{cdm} e a matriz de inércia da espaçonave.

Na formulação adotada, assume-se que o movimento do \gls{cdm} é prescrito pelo modelo de translação orbital.
Em um modelo completamente livre de forças externas, a conservação da quantidade de movimento linear implicaria um pequeno deslizamento da espaçonave ao longo da órbita em resposta ao deslocamento dos elos.

%%%%%%%%%%%%%%%

A energia cinética do $i$-ésimo elo do manipulador é composta pelas parcelas translacional e rotacional, sendo:

\begin{equation}
T_i =
\frac{1}{2} m_i\, {}^{B}\!\vec v_{G_i}^{\,T}\,{}^{B}\!\vec v_{G_i}
+
\frac{1}{2}\, {}^{B}\!\vec\omega_i^{\,T}\, I_i\, {}^{B}\!\vec\omega_i.
\label{eq_Ti}
\end{equation}

Em que ${}^{B}\!\vec v_{G_i}$ e ${}^{B}\!\vec\omega_i$ são, respectivamente, a velocidade linear no \gls{cdm} do elo $i$
e a velocidade angular do elo $i$, expressas no referencial da espaçonave $\{B\}$ e obtidas a partir do Jacobiano acoplado introduzido na Seção~\ref{sec_cinematica_diferencial_jacobiano}, e $I_i$ é o
matriz de inércia do elo $i$ em torno do seu próprio \gls{cdm}.

%%%%%%%%%%%%%%%%%%%

Do ponto de vista físico, o termo translacional em \eqref{eq_Ti} representa o movimento do \gls{cdm} do elo $i$ em relação ao corpo da espaçonave, por meio do termo quadrático na velocidade linear ${}^{B}\!\vec v_{G_i}$.
Ao agrupar as coordenadas generalizadas em um único vetor, define-se:

\begin{equation}
\vec y
=
\begin{bmatrix}
\vec\eta_B \\
\vec{\xi}
\end{bmatrix},
\qquad
\dot{\vec y}
=
\begin{bmatrix}
\dot{\vec\eta}_B \\
\dot{\vec{\xi}}
\end{bmatrix}.
\label{eq_veloc_generalizadas}
\end{equation}

Em que $\vec\eta_B$ reúne os parâmetros de atitude da espaçonave e $\vec{\xi}$ contém as coordenadas articulares do manipulador.
Note-se que a velocidade angular da espaçonave,
${}^{B}\!\vec\omega_B$,
não é, em si, uma velocidade generalizada, mas é obtida como combinação linear das velocidades generalizadas de atitude, isto é,
${}^{B}\!\vec\omega_B = \mathbf{T}_\omega(\vec\eta_B)\,\dot{\vec\eta}_B$,
conforme a relação usual entre taxas de ângulos de Euler e velocidade angular.
Com essa definição, a energia cinética total pode ser escrita como:

\begin{equation}
T =
\frac{1}{2}\,\dot{\vec y}^{\,T}\,
M(\vec y)\,
\dot{\vec y}.
\label{eq_T_quadratica_Mq1}
\end{equation}

De forma equivalente:

\begin{equation}
T
=
\frac{1}{2}
\begin{bmatrix}
\dot{\vec{\eta}}_{B}\\
\dot{\vec{\xi}}
\end{bmatrix}^{T}
\begin{bmatrix}
M_{BB}(\vec{\eta}_B,\vec{\xi}) & M_{B\xi}(\vec{\eta}_B,\vec{\xi})\\
M_{\xi B}(\vec{\eta}_B,\vec{\xi}) & M_{\xi\xi}(\vec{\eta}_B,\vec{\xi})
\end{bmatrix}
\begin{bmatrix}
\dot{\vec{\eta}}_{B}\\
\dot{\vec{\xi}}
\end{bmatrix}.
\label{eq_T_quadratica_Mq2}
\end{equation}


Em que $M(\vec y)\in\mathbb{R}^{(3+n)\times(3+n)}$ é a matriz de inércia generalizada do conjunto espaçonave--manipulador, função da configuração $(\vec\eta_B,\vec{\xi})$, construída a partir das contribuições da espaçonave e dos elos segundo as expressões \eqref{eq_TB} e \eqref{eq_Ti}.
A partir da energia cinética \eqref{eq_T_quadratica_Mq1} e das coordenadas generalizadas, e aplicando as equações de Lagrange, obtém-se a forma matricial das equações de movimento do conjunto espaçonave e manipulador:

\begin{equation}
M(\vec y)\,\ddot{\vec y} + C(\vec y,\dot{\vec y})\,\dot{\vec y} = \vec\tau.
\label{eq_eom_acoplada_compacta}
\end{equation}

O termo $C(\vec y,\dot{\vec y})\,\dot{\vec y}$ reúne as contribuições dependentes de velocidade (Coriolis e centrífugas):


\begin{equation}
\vec\tau =
\begin{bmatrix}
\vec\tau_B \\
\tau_1 \\
\tau_2 \\
\tau_3
\end{bmatrix},
\qquad
\ddot{\vec y}=
\begin{bmatrix}
\ddot{\vec\eta}_B \\
\ddot{\vec\xi}
\end{bmatrix},
\qquad
\ddot{\vec\xi}=
\begin{bmatrix}
\ddot{\xi}_1\\
\ddot{\xi}_2\\
\ddot{\xi}_3
\end{bmatrix}.
\label{eq_tau_y_ddoty}
\end{equation}

Em que $[\cdot]_i$ denota a $i$-ésima componente do vetor resultante.
Particionando a matriz de inércia e os termos dependentes de velocidade de acordo com
$\vec y=[\,\vec\eta_B^{\,T}\ \ \vec{\xi}^{\,T}\,]^{T}$, obtém-se:

\begin{equation}
M(\vec y)=
\begin{bmatrix}
M_{BB}(\vec y) & M_{B\xi}(\vec y)\\
M_{\xi B}(\vec y) & M_{\xi\xi}(\vec y)
\end{bmatrix},
\qquad
C(\vec y,\dot{\vec y})=
\begin{bmatrix}
C_{BB}(\vec y,\dot{\vec y}) & C_{B\xi}(\vec y,\dot{\vec y})\\
C_{\xi B}(\vec y,\dot{\vec y}) & C_{\xi\xi}(\vec y,\dot{\vec y})
\end{bmatrix}.
\label{eq_particionamento_M_C}
\end{equation}

Tomando, em \eqref{eq_eom_acoplada_compacta}, as linhas associadas às coordenadas articulares $\vec{\xi}$ e empregando o particionamento \eqref{eq_particionamento_M_C}, obtêm-se as equações escalares das juntas.
Em particular, ao selecionar a primeira linha do bloco associado às coordenadas articulares, obtém-se a equação escalar da primeira junta:

\begin{equation}
\tau_1
=
\bigl[M_{\xi B}(\vec y)\,\ddot{\vec\eta}_B\bigr]_1
+
\bigl[M_{\xi\xi}(\vec y)\,\ddot{\vec\xi}\bigr]_1
+
\bigl[C_{\xi B}(\vec y,\dot{\vec y})\,\dot{\vec\eta}_B\bigr]_1
+
\bigl[C_{\xi\xi}(\vec y,\dot{\vec y})\,\dot{\vec\xi}\bigr]_1.
\label{eq_eom_junta1}
\end{equation}

De modo análogo, a segunda junta é descrita pela segunda linha do mesmo bloco:

\begin{equation}
\tau_2
=
\bigl[M_{\xi B}(\vec y)\,\ddot{\vec\eta}_B\bigr]_2
+
\bigl[M_{\xi\xi}(\vec y)\,\ddot{\vec\xi}\bigr]_2
+
\bigl[C_{\xi B}(\vec y,\dot{\vec y})\,\dot{\vec\eta}_B\bigr]_2
+
\bigl[C_{\xi\xi}(\vec y,\dot{\vec y})\,\dot{\vec\xi}\bigr]_2.
\label{eq_eom_junta2}
\end{equation}

Por fim, a terceira junta é obtida ao selecionar a terceira linha correspondente:

\begin{equation}
\tau_3
=
\bigl[M_{\xi B}(\vec y)\,\ddot{\vec\eta}_B\bigr]_3
+
\bigl[M_{\xi\xi}(\vec y)\,\ddot{\vec\xi}\bigr]_3
+
\bigl[C_{\xi B}(\vec y,\dot{\vec y})\,\dot{\vec\eta}_B\bigr]_3
+
\bigl[C_{\xi\xi}(\vec y,\dot{\vec y})\,\dot{\vec\xi}\bigr]_3.
\label{eq_eom_junta3}
\end{equation}

Essas expressões exibem a contribuição dos termos inerciais e dos termos dependentes de velocidade provenientes do acoplamento entre a atitude do perseguidor e as coordenadas articulares do manipulador.


